% Lua Integration Showcase
% Demonstrates LuaLaTeX's embedded Lua engine with practical examples

\documentclass[
    language=english,
    doctype=article,
    institution=none,
]{../../omnilatex}

\RequirePackage{config/document-settings}
\usepackage{luacode}

%% ── Shared Lua preamble ──────────────────────────────
\begin{luacode*}
-- Date formatting helper
function smartdate(format_str)
    local t = os.date("*t")
    local months = {
        "January","February","March","April","May","June",
        "July","August","September","October","November","December"
    }
    if format_str == "long" then
        tex.sprint(months[t.month] .. " " .. t.day .. ", " .. t.year)
    elseif format_str == "iso" then
        tex.sprint(string.format("%04d-%02d-%02d", t.year, t.month, t.day))
    else
        tex.sprint(os.date("%d %B %Y"))
    end
end

-- Fibonacci with memoization
local fib_cache = {[0]=0, [1]=1}
function fib(n)
    if fib_cache[n] then return fib_cache[n] end
    fib_cache[n] = fib(n-1) + fib(n-2)
    return fib_cache[n]
end

-- Prime checker
function is_prime(n)
    if n < 2 then return false end
    for i = 2, math.floor(math.sqrt(n)) do
        if n % i == 0 then return false end
    end
    return true
end

-- KV store
local texdata = {}
function texdata_add(key, value)
    texdata[key] = value
end
function texdata_get(key)
    return texdata[key] or "?"
end

-- Factorial
function factorial(n)
    if n <= 1 then return 1 end
    return n * factorial(n-1)
end
\end{luacode*}

%% ── LaTeX wrappers ───────────────────────────────────
\newcommand{\smartdate}[1]{\luadirect{smartdate("#1")}}
\newcommand{\storedata}[2]{\luadirect{texdata_add("#1", "#2")}}
\newcommand{\retrieve}[1]{\luadirect{tex.sprint(texdata_get("#1"))}}

\begin{document}

\maketitle
\tableofcontents
\newpage

% ==================================================================
\section{Inline Lua Execution}
% ==================================================================

The \texttt{\textbackslash directlua} command executes Lua code and prints the
result of \texttt{tex.sprint()} calls directly into the document.

2 + 2 = \luadirect{tex.sprint(2 + 2)}

Lua can access the current date and time:
\luadirect{
    local t = os.date("*t")
    tex.sprint(string.format(
        "Year: \%d, Month: \%d, Day: \%d",
        t.year, t.month, t.day
    ))
}

% ==================================================================
\section{Lua Functions}
% ==================================================================

Define reusable Lua functions in the preamble, then call them inline.

\begin{luacode*}
function greet(name)
    tex.sprint("Hello, " .. name .. "! Welcome to LuaLaTeX.")
end
\end{luacode*}

\luadirect{greet("World")}

\luadirect{greet("OmniLaTeX User")}

% ==================================================================
\section{Table Manipulation}
% ==================================================================

Lua tables are powerful data structures. Here we create a sorted list of
programming languages:

\begin{itemize}
\luadirect{
    local langs = {"Lua", "Python", "Haskell", "Rust", "Go"}
    table.sort(langs)
    for _, lang in ipairs(langs) do
        tex.sprint("\\item " .. lang)
    end
}
\end{itemize}

\medskip

A word-count table built from a sentence:

\luadirect{
    local sentence = "the quick brown fox jumps over the lazy dog"
    local counts = {}
    for word in sentence:gmatch("\%S+") do
        counts[word] = (counts[word] or 0) + 1
    end
    tex.sprint("\\begin{tabular}{lc}")
    tex.sprint("\\toprule")
    tex.sprint("\\textbf{Word} & \\textbf{Count} \\\\")
    tex.sprint("\\midrule")
    for w, c in pairs(counts) do
        tex.sprint(w .. " & " .. c .. " \\\\")
    end
    tex.sprint("\\bottomrule")
    tex.sprint("\\end{tabular}")
}

% ==================================================================
\section{String Processing}
% ==================================================================

Lua's pattern-matching library is ideal for text transformations.

\luadirect{
    local s = "hello_world_example"
    local title = s:gsub("_", " "):gsub("(%a)([\%w]*)", function(first, rest)
        return first:upper() .. rest:lower()
    end)
    tex.sprint("Title case: " .. title)
}

\medskip

Extracting all numbers from a string:

\luadirect{
    local data = "version 3.14 has 42 improvements and 7 bug fixes"
    local nums = {}
    for n in data:gmatch("\%d+\.?\%d*") do
        nums[#nums + 1] = n
    end
    tex.sprint("Numbers found: " .. table.concat(nums, ", "))
}

% ==================================================================
\section{Mathematical Computation}
% ==================================================================

\begin{center}
\luadirect{
    tex.sprint("\\begin{tabular}{r*{6}{c}}")
    tex.sprint("\\toprule")
    tex.sprint("$n$     & 1 & 2 & 3 & 4 & 5 & 6 \\\\")
    tex.sprint("\\midrule")
    tex.sprint("$F(n)$  & " .. fib(1) .. " & " .. fib(2) .. " & "
               .. fib(3) .. " & " .. fib(4) .. " & "
               .. fib(5) .. " & " .. fib(6) .. " \\\\")
    tex.sprint("\\bottomrule")
    tex.sprint("\\end{tabular}")
}
\end{center}

\medskip

\begin{center}
\luadirect{
    tex.sprint("\\begin{tabular}{r*{6}{c}}")
    tex.sprint("\\toprule")
    tex.sprint("$n$     & 7 & 8 & 9 & 10 & 11 & 12 \\\\")
    tex.sprint("\\midrule")
    tex.sprint("$F(n)$  & " .. fib(7) .. " & " .. fib(8) .. " & "
               .. fib(9) .. " & " .. fib(10) .. " & "
               .. fib(11) .. " & " .. fib(12) .. " \\\\")
    tex.sprint("\\bottomrule")
    tex.sprint("\\end{tabular}")
}
\end{center}

% ==================================================================
\section{Custom Lua-Backed Commands}
% ==================================================================

Define \LaTeX{} commands whose implementations live entirely in Lua.

Long date: \smartdate{long} \quad ISO date: \smartdate{iso}

\medskip

A \texttt{\textbackslash luasum} command that sums numbers in a string:

\newcommand{\luasum}[1]{\luadirect{%
    local sum = 0
    for v in "#1":gmatch("-?\%d+\.?\%d*") do sum = sum + tonumber(v) end
    tex.sprint(sum)
}}

\luasum{10,20,30,40} \quad \luasum{1.5 2.5 3}

% ==================================================================
\section{Prime Number Table}
% ==================================================================

Generate a table of the first 25 prime numbers using Lua:

\begin{center}
\luadirect{
    tex.sprint("\\begin{tabular}{*{5}{c}}")
    tex.sprint("\\toprule")
    local count, n = 0, 2
    while count < 25 do
        if is_prime(n) then
            count = count + 1
            tex.sprint(n)
            if count % 5 == 0 then
                tex.sprint(" \\\\")
            elseif count < 25 then
                tex.sprint(" & ")
            end
        end
        n = n + 1
    end
    tex.sprint("\\bottomrule")
    tex.sprint("\\end{tabular}")
}
\end{center}

% ==================================================================
\section{Practical Example: CSV Table}
% ==================================================================

Parse CSV data with Lua and render it as a \LaTeX{} table with automatic
column alignment detection.

\luadirect{
    local csv = [[
Name,Role,Experience,Start Year
Alice,Engineer,5,2020
Bob,Designer,3,2022
Carol,Manager,8,2017
Dave,Engineer,6,2019
Eve,Analyst,2,2023
]]

    local rows = {}
    for line in csv:gmatch("[^\r\n]+") do
        local row = {}
        for field in line:gmatch("([^,]+)") do
            row[#row + 1] = field
        end
        if #row > 0 then rows[#rows + 1] = row end
    end

    local ncols = #rows[1]
    local colspec = "l" .. string.rep("c", ncols - 1)

    tex.sprint("\\begin{table}[htbp]")
    tex.sprint("\\centering")
    tex.sprint("\\caption{Team members generated from CSV data}")
    tex.sprint("\\label{tab:csv}")
    tex.sprint("\\begin{tabular}{" .. colspec .. "}")
    tex.sprint("\\toprule")

    for i, row in ipairs(rows) do
        if i == 1 then
            tex.sprint("\\textbf{" .. table.concat(row, "} & \\textbf{") .. "} \\\\")
            tex.sprint("\\midrule")
        else
            tex.sprint(table.concat(row, " & ") .. " \\\\")
        end
    end

    tex.sprint("\\bottomrule")
    tex.sprint("\\end{tabular}")
    tex.sprint("\\end{table}")
}

% ==================================================================
\section{TeX--Lua Interaction}
% ==================================================================

Pass data from \LaTeX{} into Lua and back:

\storedata{project}{OmniLaTeX}
\storedata{version}{2.3}
\storedata{engine}{LuaLaTeX}

Project name: \retrieve{project}, version \retrieve{version}, built with \retrieve{engine}.

% ==================================================================
\section{Factorial Table}
% ==================================================================

Compute factorials directly in Lua and display as a table:

\begin{center}
\luadirect{
    tex.sprint("\\begin{tabular}{rrc}")
    tex.sprint("\\toprule")
    tex.sprint("\\multicolumn{1}{c}{$n$} & \\multicolumn{1}{c}{$n!$} & \\multicolumn{1}{c}{Prime?} \\\\")
    tex.sprint("\\midrule")
    for n = 1, 15 do
        local f = factorial(n)
        tex.sprint(n .. " & " .. f .. " & ")
        if is_prime(n) then
            tex.sprint("\\checkmark")
        else
            tex.sprint("---")
        end
        tex.sprint(" \\\\")
    end
    tex.sprint("\\bottomrule")
    tex.sprint("\\end{tabular}")
}
\end{center}

% ==================================================================
\section{Random Data Generation}
% ==================================================================

Lua's \texttt{math.random} can generate sample data for figures and tables:

\begin{center}
\luadirect{
    tex.sprint("\\begin{tikzpicture}[scale=0.8]")
    tex.sprint("\\draw[->] (0,0) -- (10.5,0) node[right] {$x$};")
    tex.sprint("\\draw[->] (0,0) -- (0,6.5) node[above] {$y$};")
    math.randomseed(42)
    for i = 1, 30 do
        local x = math.random() * 10
        local y = 2 + 0.3 * x + math.random() * 2
        tex.sprint(string.format(
            "\\fill[blue!60] (%.2f,%.2f) circle (0.12);", x, y
        ))
    end
    tex.sprint("\\end{tikzpicture}")
}
\end{center}

% ==================================================================
\section{File I/O: Read External Data}
% ==================================================================

Lua can read files from disk during compilation.  This example reads
the \texttt{polyglossia} settings from \texttt{lib/language/omnilatex-i18n.sty}
and displays the first few language declarations.

\luadirect{
    local f = io.open("lib/language/omnilatex-i18n.sty", "r")
    if f then
        tex.sprint("\\begin{itemize}")
        local count = 0
        for line in f:lines() do
            local lang = line:match("setdefaultlanguage{(.+)}")
                 or line:match("setotherlanguages{(.+)}")
            if lang then
                count = count + 1
                tex.sprint("\\item \\texttt{" .. lang .. "}")
                if count >= 5 then break end
            end
        end
        f:close()
        tex.sprint("\\item \\dots")
        tex.sprint("\\end{itemize}")
    else
        tex.sprint("(File not found --- running in sandboxed CI)")
    end
}

\end{document}
